%==============================================================================
\section{Data Loading --- Nạp dữ liệu vào dự án học máy}
%==============================================================================

\begin{ynghia}[title={Nguồn}]
\tsurl{https://www.tutorialspoint.com/machine_learning/machine_learning_data_loading.htm}.
\end{ynghia}

\begin{dongco}[title={Bước đầu tiên của mọi dự án ML}]
Khi bắt đầu dự án ML, thứ cần nhất là \textbf{dữ liệu} phải được nạp vào hệ thống.
\end{dongco}

\begin{dinhnghia}[title={Data loading}]
Trong ML, data loading là quá trình nhập/đọc dữ liệu từ nguồn bên ngoài và chuyển sang định dạng thuật toán dùng được.
Sau đó dữ liệu được tiền xử lý (loại bất thường, giá trị thiếu, outlier).
Khi đã tiền xử lý, dữ liệu được chia train/test để huấn luyện và đánh giá mô hình.
\end{dinhnghia}

\begin{ynghia}[title={Nguồn và định dạng}]
Dữ liệu có thể từ CSV, database, web API, cloud storage, \ldots
Định dạng phổ biến nhất cho dự án ML là \textbf{CSV} (Comma Separated Values).
\end{ynghia}

\subsection{Lưu ý khi nạp dữ liệu CSV}

\begin{dinhnghia}[title={CSV}]
CSV là định dạng văn bản thuần lưu dữ liệu dạng bảng: mỗi hàng là một bản ghi, mỗi cột là một trường/thuộc tính.
Đơn giản, nhẹ, dễ đọc/xử lý bằng Python, R, Java.
Trong Python có nhiều cách nạp CSV vào dự án ML; trước khi nạp cần lưu ý các thành phần sau.
\end{dinhnghia}

\subsubsection{File Header (hàng tiêu đề)}

\begin{phuongphap}[title={Vai trò và lưu ý}]
Hàng đầu tiên của CSV thường chứa tên cột.
\begin{itemize}
  \item \textbf{Consistency}: số cột và tên cột phải nhất quán trên mọi hàng.
  \item \textbf{Meaningful names}: tên cột có nghĩa, mô tả; tránh \texttt{column1}, \texttt{column2}, \ldots
  \item \textbf{Case sensitivity}: tên cột có thể phân biệt hoa/thường tùy thư viện.
  \item \textbf{Special characters}: tránh khoảng trắng, dấu phẩy, ngoặc kép trong tên cột; dùng underscore hoặc camelCase.
  \item \textbf{Missing header}: nếu không có header, cần chỉ định tên cột thủ công hoặc tài liệu kèm theo.
  \item \textbf{Encoding}: encoding của header phải tương thích với công cụ đọc file.
\end{itemize}
\end{phuongphap}

\subsubsection{Comments (dòng chú thích)}

\begin{ynghia}[title={Dòng comment}]
Dòng tùy chọn bắt đầu bằng ký tự như \texttt{\#} hoặc \texttt{//}, thường bị bỏ qua khi đọc CSV.
Không thường là dữ liệu huấn luyện, nhưng nếu có cần xử lý đúng khi nạp.
\end{ynghia}

\begin{itemize}
  \item Biết \textbf{comment marker} đang dùng.
  \item Comment nên ở \textbf{dòng riêng}, không trộn với dữ liệu.
  \item Dùng marker \textbf{nhất quán} trong cả file.
  \item Hiểu thư viện có bỏ qua comment mặc định hay cần tham số.
  \item Nếu comment chứa thông tin quan trọng, có thể cần xử lý tách khỏi dữ liệu.
\end{itemize}

\subsubsection{Delimiter (ký tự phân tách)}

\begin{luuy}[title={Ảnh hưởng tới mô hình}]
Delimiter ảnh hưởng đáng kể độ chính xác và hiệu năng khi nạp dữ liệu.
\begin{itemize}
  \item Chọn delimiter phù hợp (ví dụ giá trị có dấu phẩy bên trong $\Rightarrow$ có thể dùng tab hoặc semicolon thay vì comma).
  \item \textbf{Consistency} trong cả file.
  \item \textbf{Encoding} tương thích với delimiter (ví dụ delimiter không-ASCII với UTF-8).
  \item Một số thư viện ML yêu cầu delimiter cụ thể --- xem tài liệu công cụ.
\end{itemize}
\end{luuy}

\subsubsection{Quotes (dấu ngoặc kép)}

\begin{phuongphap}[title={Lưu ý khi dùng quotes}]
\begin{itemize}
  \item Ký tự quote (\texttt{"} phổ biến nhất) nhất quán trong file.
  \item Giá trị có delimiter hoặc xuống dòng có thể được bọc trong quotes.
  \item Quote bên trong giá trị phải \textbf{escape} (thường nhân đôi quote), ví dụ: \texttt{"John ""the Hammer"" Smith"}.
  \item Cách dùng quote có thể khác nhau tùy công cụ tạo CSV.
  \item Encoding ảnh hưởng cách đọc quotes.
\end{itemize}
\end{phuongphap}

\subsection{Các cách nạp file CSV trong Python}

\begin{ynghia}[title={Nhiệm vụ then chốt}]
Nạp đúng dữ liệu là bước quan trọng nhất; CSV có nhiều biến thể và độ khó parse khác nhau.
\end{ynghia}

\subsubsection{Dùng module CSV (built-in)}

\begin{vidu}[title={csv.reader}]
\begin{lstlisting}
import csv
with open('mydata.csv', 'r') as file:
   reader = csv.reader(file)
   for row in reader:
      print(row)
\end{lstlisting}
Đọc \texttt{mydata.csv} và in từng hàng.
\end{vidu}

\subsubsection{Dùng thư viện Pandas}

\begin{vidu}[title={pandas.read\_csv}]
\begin{lstlisting}
import pandas as pd

data = pd.read_csv('mydata.csv')
\end{lstlisting}
Nạp CSV vào DataFrame \texttt{data} --- tiện cho thao tác và biến đổi dữ liệu.
\end{vidu}

\subsubsection{Dùng thư viện NumPy}

\begin{vidu}[title={numpy.genfromtxt}]
\begin{lstlisting}
import numpy as np

data = np.genfromtxt('mydata.csv', delimiter=',')
\end{lstlisting}
Nạp CSV vào mảng NumPy \texttt{data}.
\end{vidu}

\subsubsection{Dùng thư viện SciPy}

\begin{vidu}[title={scipy loadtxt}]
\begin{lstlisting}
import numpy as np

from scipy import loadtxt
data = loadtxt('mydata.csv', delimiter=',')
\end{lstlisting}
\end{vidu}

\subsubsection{Dùng thư viện scikit-learn}

\begin{vidu}[title={load\_iris}]
\begin{lstlisting}
from sklearn.datasets import load_iris

data = load_iris().data
\end{lstlisting}
Nạp dataset Iris có sẵn trong sklearn vào mảng NumPy \texttt{data}.
\end{vidu}

\begin{luuy}[title={Đối chiếu}]
Các bước tiền xử lý, hiểu dữ liệu và chia train/test: bài \texttt{19--23}; train/validation/test: Mục~\ref{sec:train-val-test}.
\end{luuy}
